\chapter{Architecture et Open Supply unifié}\label{ch:architecture}

\section{Quatre classes}

Quatre classes portent le mécanisme (\cref{fig:uml}). Les noms sont ceux du code à la révision
validée ; les autres attributs ne sont pas représentés.

\begin{itemize}
  \item \jv{CommandeAgent} : une commande client ou un ordre autonome \jv{REAPPRO\_n}. Elle porte
    son type (MTS, MTO ou ETO), sa quantité, son échéance et son statut.
  \item \jv{FicheMatiere} : une matière. Elle porte le stock physique (\jv{stockDisponible}), les
    réservations, les réceptions attendues et, optionnellement, un buffer DDMRP.
  \item \jv{DDMRPBuffer} : les paramètres et l'état du buffer d'une matière (ADU, lead time, zones,
    NFP, Qualified Demand décomposée, statut). Il ne contient pas de logique : les calculs sont
    dans les fonctions de la classe principale.
  \item \jv{ReceptionAttendue} : une quantité commandée et non encore reçue, avec son échéance,
    son origine, son statut et ses allocations à des commandes.
\end{itemize}

\begin{figure}[H]
  \centering
  \begin{tikzpicture}[font=\small,
    cls/.style={draw=ink, rounded corners=2pt, fill=white, align=left, inner sep=3pt, text width=3.9cm, font=\scriptsize},
    rel/.style={-{Stealth[length=2.2mm]}, ink, thick},
    lbl/.style={font=\scriptsize, fill=white, inner sep=1pt}]
    \node[cls] (cmd) at (0,0) {\textbf{CommandeAgent}\\[2pt]
      \texttt{idCommande}\\\texttt{typeCommande}\\(MTS / MTO / ETO)\\\texttt{qte}\\\texttt{tPromise}\\\texttt{statut}};
    \node[cls] (fm) at (5.6,0) {\textbf{FicheMatiere}\\[2pt]
      \texttt{idMatiere}\\\texttt{stockDisponible}\\\texttt{reservations}\\\texttt{receptionsAttendues}\\\texttt{consommeParCommande}};
    \node[cls] (rec) at (11.2,0) {\textbf{ReceptionAttendue}\\[2pt]
      \texttt{quantite}\\\texttt{dateArriveePrevue}\\\texttt{origine}\\(S\_Q, DDMRP, MTS, MTO, ETO)\\\texttt{statut}\\(OUVERTE / RECUE)\\\texttt{allocations}};
    \node[cls] (db) at (5.6,-6.2) {\textbf{DDMRPBuffer}\\[2pt]
      \texttt{actif}\\\texttt{adu}\\\texttt{leadTimeHeures}\\\texttt{zoneRouge}, \texttt{zoneJaune}\\\texttt{zoneVerte}\\\texttt{topOfRed}, \texttt{topOfYellow}\\\texttt{topOfGreen}\\\texttt{openSupply}\\\texttt{qualifiedDemand}\\\texttt{netFlowPosition}\\\texttt{statut}, \texttt{priorite}};
    \draw[rel] (cmd.east) -- node[lbl, above] {nomenclature} (fm.west);
    \draw[rel] (fm.east) -- node[lbl, above] {0..*} (rec.west);
    \draw[rel] (fm.south) -- node[lbl, right] {0..1} (db.north);
    \draw[rel, dashed] (cmd.south) |- node[lbl, pos=0.25, below] {allocation} (rec.south);
  \end{tikzpicture}
  \caption{Vue simplifiée des classes : la fiche matière possède le stock, les réservations,
    les réceptions attendues et, optionnellement, le buffer DDMRP.}
  \label{fig:uml}
\end{figure}


La fiche matière est la source de vérité pour le stock. Une réservation ne modifie pas ce stock ;
elle réduit la disponibilité opérationnelle. Le buffer ne possède aucun stock : il décrit une
situation et en déduit une décision. Une réception n'est jamais un stock : elle le devient lorsqu'elle
est créditée.

\section{Open Supply : une représentation unique}

Toute matière commandée mais non encore reçue est une \jv{ReceptionAttendue} dans la fiche
matière. L'ensemble de ces réceptions forme l'Open Supply de la matière. Cette représentation
unique répond à un problème précis. Auparavant, l'approvisionnement lancé par une commande MTO
était crédité directement au moment de la finalisation, sans être visible pour les autres
décisions. Une politique DDMRP, qui doit savoir ce qui est déjà commandé, ne pouvait pas le voir.

Le cycle de vie d'une réception est représenté sur \cref{fig:open-supply}. Elle est créée, reste
ouverte jusqu'à son échéance, est créditée une seule fois, puis retirée de l'Open Supply.
Deux fonctions portent ce cycle. \jv{creerReceptionOpenSupply} est le seul point de création.
\jv{traiterReceptionsEchues} est le seul point de crédit du stock : il ajoute la quantité au stock
physique et retire la réception de la liste dans la même opération. Une réception ne peut donc pas
être créditée deux fois.

\begin{figure}[H]
  \centering
  \begin{tikzpicture}[font=\small, >=Stealth, node distance=1.5cm,
    st/.style={draw=ink, rounded corners=3pt, minimum width=5.2cm, minimum height=0.95cm, align=center, fill=white},
    tr/.style={-{Stealth[length=2.2mm]}, thick},
    lb/.style={font=\scriptsize, align=left, text width=4.6cm}]
    \node[st, fill=ink!8] (cr) {Création\\\texttt{creerReceptionOpenSupply}};
    \node[st, fill=osc!22, below=of cr] (op) {OUVERTE : compte dans Open Supply};
    \node[st, fill=ddmrpgrn!22, below=of op] (rc) {RECUE : crédit physique\\\texttt{traiterReceptionsEchues}};
    \node[st, fill=ink!8, below=of rc] (sup) {Retirée de la liste : Open Supply réduit};
    \draw[tr] (cr) -- node[lb, right, xshift=2.7cm] {ajout de la réception} (op);
    \draw[tr] (op) -- node[lb, right, xshift=2.7cm] {échéance atteinte} (rc);
    \draw[tr] (rc) -- node[lb, right, xshift=2.7cm] {retrait immédiat} (sup);
    \draw[tr, dashed] (op.west) -- ++(-0.9,0) |- node[lb, left, pos=0.25, xshift=-0.1cm, align=right, text width=3.2cm] {allocation d'une part à une commande}
      ($(op.west)+(-0.9,0)$);
  \end{tikzpicture}
  \caption{Cycle de vie d'une réception : elle est créée, reste ouverte jusqu'à l'échéance,
    est créditée une seule fois et disparaît de l'Open Supply.}
  \label{fig:open-supply}
\end{figure}


\section{Allocation et réservation}

Une commande peut recevoir une allocation sur une réception ouverte. Cette allocation signifie
« cette commande attend cette part de la réception ». Elle sert à la coordination du workflow MTO :
une commande n'ouvre pas une seconde réception pour une quantité déjà couverte. L'allocation ne
réduit pas la demande qualifiée du DDMRP. Elle n'est pas non plus une réservation physique : la
réservation n'a lieu qu'une fois la matière effectivement créditée, et elle est décrite au
\cref{ch:mto2}.

\section{Origine des réceptions}

L'origine trace la décision qui a créé la réception. Cinq valeurs sont possibles. \texttt{S\_Q}
désigne la politique (s,Q) historique, pour les matières non gérées par DDMRP. \texttt{DDMRP}
désigne une décision du buffer. \texttt{MTS}, \texttt{MTO} et \texttt{ETO} désignent le type réel de
la commande ayant déclenché l'approvisionnement. Les ordres REAPPRO sont des commandes MTS : leurs
réceptions portent l'origine \texttt{MTS}. L'origine sert à la traçabilité, et aucune décision ne la lit.
